% Transcription — fonds Alexandre Grothendieck, shelfmark 19, batch 2 (pages 21-40).
% Established from the facsimile archives/batches/19/batch-02.pdf (a mirror of
% https://grothendieck.umontpellier.fr/19.pdf), per the skill
% transcribe-grothendieck.
%
% Pages 29 and 34 are blank versos — absent.
%
% The batch holds the end of one piece and the whole of another. Pages 21-27
% finish "Cohomologie relative", begun on page 13 and transcribed in batch 1:
% when the S-derived functors of Gamma agree with the ordinary ones, the
% characterisation of the S-injectives as the direct factors of the T(M), the
% S-resolutions and the effacing morphisms, the canonical resolution C^*(A),
% and three examples — restriction of scalars along a subring V of U, its
% adjoint on the other side, and Godement's C^0 for sheaves of O-modules.
%
% Page 28 is the wrapper of the second piece and names it: "Remarques —
% SGA 4 IX", with a pencilled tally "11 p" that matches pages 30-40 exactly.
% What follows is not a draft of the exposé but notes towards it, in a much
% more compressed hand: the plan of nos 1-3 with eight equivalent conditions
% for constructibility (page 30), a purity proposition (31), the rest of the
% plan and two "Leitfaden" sketches (32), a descent statement for proper
% morphisms with a criterion on the geometric fibres (33), then XII 6.2 —
% the five-term ladder comparing two cohomological functors and the
% bookkeeping of the conditions I_n, J_n (35, 37), the proposition on a class
% M of morphisms of locally noetherian preschemes closed under the four
% conditions a), b), c), d), from which every projective morphism lies in M
% and, via Chow's lemma, every proper one (36, 38), a devissage for
% constructible l-torsion sheaves (39), and a crossed-out page on R h^!
% (40). The sheets are not in reading order: page 37 continues on page 39,
% and page 38 states the proposition whose corollary page 36 draws.
%
% This second piece is written fast and small, and a large part of its
% connecting prose could not be read: the \ill{} marks are frequent and they
% are literal. What is legible is the mathematical skeleton — the statements,
% the conditions, the diagrams — and that is what is given here.
%
% The dating is the inventory's own, for the whole shelfmark; nothing on these
% pages carries a date.
%
% Pass: Opus 5 (claude-opus-5), 2026-08-22 — first pass, unchecked against the pages by a human.
\documentclass[11pt,a4paper]{article}
% Le préambule commun à toutes les transcriptions du fonds Grothendieck.
%
% Il tient en une idée : **l'incertitude se compose, elle ne se résout pas.**
% Une transcription qui lisse un mot douteux a détruit la seule chose pour
% laquelle on la faisait — un lecteur doit pouvoir savoir, à chaque endroit,
% ce qui était sur le feuillet et ce qui a été ajouté. Les six macros qui
% suivent sont donc l'appareil critique, et non de la décoration.
%
% Les mêmes commandes sont reconnues par `scripts/render.mjs`, qui produit la
% vue de lecture du site. Une macro ajoutée ici doit l'être aussi là-bas,
% faute de quoi le rendu échouera bruyamment — ce qui est voulu : un rendu
% silencieusement faux est pire qu'un rendu absent.

\ProvidesPackage{grothendieck}[2026/08/08 Transcriptions du fonds Grothendieck]

\usepackage{amsmath,amssymb,amsthm}
\usepackage{mathrsfs}
\usepackage{stmaryrd}
% Les diagrammes commutatifs sont partout dans ce fonds : c'est la langue
% même de la géométrie algébrique de Grothendieck, et une transcription qui
% les rendrait en prose ne transcrirait rien.
\usepackage{tikz-cd}
% Français par défaut — c'est la langue des feuillets — mais l'anglais est
% chargé aussi : la traduction et le résumé sont des documents anglais, et la
% typographie française (espaces avant les deux-points) y serait une faute.
% Ces documents basculent par \selectlanguage{english} en tête de corps.
\usepackage[english,french]{babel}
\usepackage{fontspec}
\usepackage{geometry}
\geometry{a4paper,margin=25mm,marginparwidth=28mm}
\usepackage{marginnote}
\usepackage{xcolor}
% ulem, not soul. `soul`'s \st and \ul are built on a character-by-character
% scan that XeLaTeX's font machinery breaks: the document fails with an
% undefined control sequence rather than degrading. ulem does the same two jobs
% and is Unicode-engine safe. `normalem` stops it hijacking \emph, which the
% transcriptions use for Grothendieck's own emphasis.
\usepackage[normalem]{ulem}
\usepackage{hyperref}
\hypersetup{colorlinks=true,linkcolor=black,urlcolor=black,pdfborder={0 0 0}}

% --- Métadonnées du lot -----------------------------------------------------
% Reprises telles quelles par le script de rendu, qui les affiche en tête de
% la vue de lecture. Elles fixent ce que le fichier prétend couvrir : sans
% elles, une transcription flotte sans point d'attache dans un fonds de seize
% mille feuillets.

\newcommand{\folder}[1]{\gdef\grothFolder{#1}}
\newcommand{\batch}[1]{\gdef\grothBatch{#1}}
\newcommand{\leaves}[2]{\gdef\grothFirst{#1}\gdef\grothLast{#2}}
\newcommand{\foldertitle}[1]{\gdef\grothFoldertitle{#1}}
\gdef\grothFolder{?}\gdef\grothBatch{?}\gdef\grothFirst{?}\gdef\grothLast{?}\gdef\grothFoldertitle{}

% --- Le résumé --------------------------------------------------------------
%
% Il ouvre la lecture modernisée, sous le titre « Résumé », et il est composé
% en petit. Ce n'est pas un ornement : c'est le seul passage du document qui
% ne soit pas des mathématiques, et un lecteur doit pouvoir le reconnaître
% avant de l'avoir lu — puis le sauter s'il connaît déjà le sujet, ou s'y
% arrêter s'il ne le connaît pas. Le filet qui le suit marque la reprise du
% corps.

\newenvironment{resume}
  {\small\setlength{\parskip}{0.45em}}
  {\par\medskip\hrule\medskip}

% --- Mots-clés ---------------------------------------------------------------
%
% En anglais, en clôture du résumé de la lecture modernisée : le vocabulaire
% moderne sous lequel chercher ce que les feuillets construisent. Le manifeste
% du site (`npm run manifest`) les extrait pour étiqueter la cote entière —
% cette ligne est la source unique des tags, il n'y a pas de fichier à côté.
\newcommand{\keywords}[1]{\par\smallskip\noindent
  {\footnotesize\textsc{Keywords} --- \textit{#1}}\par}

% --- L'appareil critique ----------------------------------------------------

% Le numéro de feuillet, en marge. C'est l'ancre qui permet de régler un
% désaccord sur une formule en regardant *un* feuillet plutôt qu'une chemise
% de six cents. Le site s'en sert aussi pour tourner les pages du fac-similé
% quand on fait défiler la transcription.
\newcommand{\leaf}[1]{%
  \marginnote{\footnotesize\color{gray!70}\textsf{#1}}%
  \label{leaf:#1}\ignorespaces}

% Illisible : on ne devine pas. Un mot inventé qui se lit comme les autres est
% le pire résultat possible.
\newcommand{\ill}{\textcolor{red!60!black}{[\,\ldots\,]}}

% Lecture douteuse : le mot est donné, et signalé comme incertain.
\newcommand{\uncertain}[1]{\uline{#1}}

% Ajout de l'éditeur — une lettre manquante, un mot sous-entendu.
\newcommand{\add}[1]{\textcolor{blue!50!black}{[#1]}}

% Biffé par Grothendieck lui-même. Ce qu'il a rayé fait partie du document :
% une rature dit souvent où la pensée a changé de direction.
\newcommand{\struck}[1]{\sout{#1}}

% Note du transcripteur, sur l'état matériel du feuillet.
\newcommand{\note}[1]{\par\smallskip{\small\color{gray!60!black}\textit{[#1]}}\par\smallskip}

% Note marginale de Grothendieck, distincte du corps du texte.
\newcommand{\marginal}[1]{\par\smallskip\noindent\hspace{1em}%
  {\small\color{gray!60!black}#1}\par\smallskip}

% --- Titre ------------------------------------------------------------------

\AtBeginDocument{%
  \begin{center}
    {\large\bfseries Fonds Alexandre Grothendieck --- cote n\textsuperscript{o}~\grothFolder}\\[2pt]
    {\itshape\grothFoldertitle}\\[4pt]
    {\small feuillets \grothFirst--\grothLast\ (lot \grothBatch)}\\[2pt]
    {\footnotesize\color{gray!60!black}Transcription établie d'après le fac-similé numérisé
     par l'Université de Montpellier.\\
     Les lectures incertaines sont soulignées, les illisibles notées [\,\ldots\,],
     les ajouts entre crochets.}
  \end{center}
  \vspace{1em}}

\endinput


\folder{19}
\batch{2}
\pages{21}{40}
\dating{[à partir de 1958-1973]}
\watermark{Édition de démonstration}
\foldertitle{Topos : notes manuscrites (s.d.)}

\begin{document}

\section*{Cohomologie relative (suite, pages 21 à 27)}

\page{21}
\struck{\ill{} $A \in C$. La construction}

Dans le cas où $C$ est une catégorie abélienne \uncertain{ayant} dit-on
suffisamment d'injectifs, la question se pose de savoir quand les foncteurs
dérivés \struck{du} \add{usuels} \uncertain{au} foncteur $\Gamma : C \to D$
coïncident avec les foncteurs dérivés usuels. \struck{Voici des}
\add{Voici les} conditions suffisantes raisonnables.

On veut donc en avoir que $R_S^{\bullet}\Gamma$ soit un foncteur
cohomologique. \struck{Il faut que $\varphi$ soit injectif, i.e.
$A \to TS(A)$} \ill{}
\note{une demi-page est ici récrite trois fois et biffée d'un seul trait ;
seul le résultat, $R_S^{0}\Gamma = \Gamma = R^{0}\Gamma$, se lit}

Je veux dire \ill{} à partir de la \ill{} d'effacer la \ill{} des
$R_S^{p}\Gamma$ \ill{} l'opérateur cobord $\delta$ défini par \ill{} en suite
courte $0 \to A' \to A \to A'' \to 0$. C'est donc le cas si
$\Gamma\bigl( C^{\bullet}(A) \bigr)$ est un foncteur exact \ill{}.

\begin{enumerate}
\item[(i)] $A \to C^{0}(A)$ est injectif ;
\item[(ii)] $S$ et $T$ \add{donc $C^{0}$, donc $\Gamma C^{\bullet}$} sont des
  foncteurs exacts ;
\item[(iii)] si $0 \to A' \to A \to A'' \to 0$ est une suite courte de $C$,
  alors la suite
  \[
    0 \longrightarrow \Gamma\bigl( C^{0}(A') \bigr) \longrightarrow
    \Gamma\bigl( C^{0}(A) \bigr) \longrightarrow
    \Gamma\bigl( C^{0}(A'') \bigr) \longrightarrow 0
  \]
  est exacte.
\end{enumerate}

[Il suffit que le foncteur $\Gamma T$ soit exact.] Alors
$R_S^{\bullet}\Gamma$ est bien un foncteur cohomologique ; en particulier, en
degré $0$, il coïncide avec $R^{0}\Gamma$, \ill{} on va supposer que
\struck{$A \to C^{0}(A)$ injectif} $\Gamma$ est exact à gauche.

\page{22}
\textbf{Exemples.} (i) \ill{} $M = T(N)$, car alors la suite (6) s'identifie à
\[
  0 \longrightarrow \mathrm{Hom}(\widetilde{A''}, \widetilde{N})
    \longrightarrow \mathrm{Hom}(\widetilde{A}, \widetilde{N})
    \longrightarrow \mathrm{Hom}(\widetilde{A'}, \widetilde{N})
    \longrightarrow 0 .
\]
D'ailleurs, si (6) est telle que, pour \struck{tout} $M = T(N)$, la suite (7)
soit exacte, alors (6) est $\sim$-exacte.

\struck{Si $M$ est injectif, alors} \ill{}
\struck{$M \to T(M) \to M^{*} \to 0$, $M^{*} = \mathrm{coker}\, i_M$}
\add{Supposons le foncteur $\sim$ exact et « fidèle »}

(ii) Soit $X \in C$ ; la suite
\begin{equation*}
  S_X : \quad 0 \longrightarrow X \xrightarrow{\ i_X\ } T(X)
    \longrightarrow X^{*} \longrightarrow 0
  \tag{8}
\end{equation*}
($X^{*} = \mathrm{coker}\, i_X$) est $\sim$-exacte. En effet, elle est
transformée en la suite exacte
\[
  S(S_X) : \quad 0 \longrightarrow \widetilde{X} \longrightarrow
  \widetilde{T(X)} \longrightarrow \widetilde{X^{*}} \longrightarrow 0 ,
\]
qui est \struck{scindée} \add{canoniquement scindée} en vertu de (4).

Supposons alors que \struck{pour toute} \add{$M \in C$ soit tel que, pour
toute} suite $\sim$-exacte de cette forme, la suite (7) soit exacte ; je dis
que $M$ est $\sim$-injectif. Cela sera vrai en particulier si $X = M$ pour
$0 \to M \to T(M) \to M^{*} \to 0$, donc $M$ est facteur direct de $T(M)$ ; ce
dernier étant $\sim$-injectif, il en est donc de même de $M$. Ceci montre
ceci : les objets injectifs sont exactement les facteurs directs des objets
$T(M)$.

\page{23}
\note{le haut de la page — la reprise de l'argument précédent — est biffé en
croix ; il se lit encore et il est donné ici}

\struck{celles pour lesquelles il existe une suite $\sim$-exacte de suites du
type $S(X)$. Pour le voir, on écrit, si $0 \to A' \to A \to A'' \to 0$ est
$\sim$-exacte,}

\marginal{Dorénavant $\sim$ est exact et « fidèle »}

Une \emph{$S$-résolution} d'un $A \in C$ est une résolution $X$ de $A$, munie
d'une opération d'homotopie \ill{} dans le complexe augmenté
\[
  0 \longrightarrow \widetilde{A} \longrightarrow \widetilde{X}^{0}
    \longrightarrow \widetilde{X}^{1} \longrightarrow \widetilde{X}^{2}
    \longrightarrow \cdots
\]
i.e.\ de projections \add{splittings} dans les suites exactes
\[
  0 \longrightarrow \widetilde{Z}^{i}(X^{\bullet}) \longrightarrow
  \widetilde{X}^{i} \longrightarrow \widetilde{Z}^{i+1}(X^{\bullet})
  \longrightarrow 0 \qquad (i \geqslant 0).
\]

\struck{Proposition. \ill{} $X$ \ill{} projectifs}

Un \struck{épimorphisme} morphisme $\eta : X \to Y$ de $C$ est dit
\emph{effaçant} si, pour tout $A \in C$, l'image de
$\mathrm{Hom}(\widetilde{A}, \widetilde{X})$ dans
$\mathrm{Hom}(\widetilde{A}, \widetilde{Y})$ par
$u \mapsto \widetilde{\eta}\,\widetilde{u}$ est contenue dans
l'\uncertain{image} de $\mathrm{Hom}(A, Y)$, i.e.\ si l'on a un diagramme
commutatif

\page{24}

\begin{tikzcd}[column sep=large]
\mathrm{Hom}(\widetilde{A}, \widetilde{X}) \arrow[r, "u \mapsto \widetilde{\eta}\widetilde{u}"] & \mathrm{Hom}(\widetilde{A}, \widetilde{Y}) \\
\mathrm{Hom}(A, T(X)) \arrow[u, "\approx"] \arrow[r, "u \mapsto T(\eta)u"] \arrow[dr, "\rho_A"'] & \mathrm{Hom}(A, T(Y)) \arrow[u, "\approx"'] \\
 & \mathrm{Hom}(A, Y) \arrow[u]
\end{tikzcd}

où $\rho_A$ est \ill{} unique. Cela revient à dire que l'application
$T(\eta) : T(X) \to T(Y)$ se factorise en
$T(X) \to Y \xrightarrow{\ i_Y\ } T(Y)$.

\struck{On voit que $\eta$ efface le foncteur $T(X)/\varepsilon_X(X)$.}
\struck{Démonstration, que} \struck{Exemple}
Si $\eta$ est un morphisme quelconque, il en est de même de $\lambda\eta$ pour
tout $\lambda \in \mathrm{Hom}(Y, \ill{})$. Un morphisme
$i_M : M \to T(M)$ est toujours effaçant.

\ill{} \emph{complexe droit} sur un objet $A \in C$ : un complexe \ill{} $X$
sur $A$, $A \to X$, muni de morphismes $t : T(A) \to X^{0}$,
$T(X^{i}) \to X^{i+1}$ ($i \geqslant 0$) rendant commutatifs les diagrammes

\begin{tikzcd}
A \arrow[r, "i"] \arrow[dr, "\varepsilon"'] & T(A) \arrow[d, "t"] \\
 & X^{0}
\end{tikzcd}

\begin{tikzcd}
X^{i} \arrow[r, "i_{X^{i}}"] \arrow[dr, "d^{i}"'] & T(X^{i}) \arrow[d, "t"] \\
 & X^{i+1}
\end{tikzcd}

\struck{Soit plus} Soit donc $A \to X$ une $S$-résolution de $A$, $B \to Y$ un
$t$-complexe droit sur $B$, et $u : A \to B$

\page{25}
un morphisme ; on \ill{} de façon canonique un homomorphisme $\bar{u}$ de $X$
dans $Y$ compatible avec les augmentations. Si on a déjà défini
$\bar{u}^{n} : X^{n} \to Y^{n}$ pour $n \leqslant m$ ($n \geqslant -1$), on
\ill{} comme suit : on a $s^{n+1} \in \mathrm{Hom}(X^{n+1}, \widetilde{X}^{n})$,
il définit $s'^{n+1} \in \mathrm{Hom}\bigl( X^{n+1}, T(Y^{n}) \bigr)$, donc un
morphisme $T(\bar{u}^{n}) \circ s'^{n+1} : X^{n+1} \to T(Y^{n})$, d'où,
\struck{en vertu de la} \add{par composition avec}
$t^{n} : T(Y^{n}) \to Y^{n+1}$, un morphisme
$t^{n}\, T(\bar{u}^{n})\, s'^{n+1} : X^{n+1} \to Y^{n+1}$. On pose
\[
  \bar{u}^{n+1} \;=\; t^{n}\, T(\bar{u}^{n})\, s'^{n+1} ,
\]

\begin{tikzcd}
X^{n} \arrow[r, "d^{n}"] \arrow[d, "\bar{u}^{n}"'] & X^{n+1} \arrow[d, "\bar{u}^{n+1}"] \\
Y^{n} \arrow[r, "d^{n}"] & Y^{n+1}
\end{tikzcd}

\struck{On a $\bar{u}^{n+1} d^{n} = t^{n} T(\bar{u}^{n}) s'^{n+1} d^{n}$}
On vérifie que ça marche \ill{}.

Résolution canonique $C^{\bullet}(A)$, \ill{} on va \ill{}
\[
  C^{n+1}(A) \;=\; T\bigl( C^{n}(A) / \mathrm{Im}\, C^{n-1}(A) \bigr)
  \qquad \text{(avec par convention } C^{-1}(M) = A\text{)} .
\]
\struck{$C^{\bullet}$} $C(A)$ : le foncteur \ill{} est une $t$-résolution.
C'est un foncteur en $A$ ; l'homomorphisme $C^{\bullet}(A) \to C^{\bullet}(B)$
associé à un homomorphisme $A \to B$ est en fait celui défini par \ill{} et la
$t$-structure de $C^{\bullet}(B)$.

\page{26}
\struck{Structure multiplicative}

\textbf{Exemples.} 1) Soient \struck{$U$} un anneau, $V$ un sous-anneau ;
considérons le foncteur
\[
  C = C^{U} \longrightarrow \widetilde{C} = C^{V}
  \qquad \text{restriction des scalaires.}
\]
C'est un foncteur exact, et les applications
$\mathrm{Hom}(A,B) \to \mathrm{Hom}(\widetilde{A}, \widetilde{B})$ sont
injectives. De plus, on a
\[
  \mathrm{Hom}_{V}(\widetilde{A}, \widetilde{B})
  \;=\; \mathrm{Hom}_{V}(A, B)
  \;\xrightarrow{\ \sim\ }\;
  \struck{\mathrm{Hom}_{U}(U \otimes_{V} A, B)}\;
  \mathrm{Hom}_{U}\bigl( A, \mathrm{Hom}_{V}(U, B) \bigr) ,
\]
qui est \struck{un isomorphisme} un \ill{} de $V$-homomorphismes naturel.

\struck{$B \xrightarrow{\sim} \mathrm{Hom}_{U}(U,B) \to \mathrm{Hom}_{V}(U,B)$}
\[
  e_B : \mathrm{Hom}_{V}(U, B) \longrightarrow B
\]
défini par $e_B(f) = f(1)$. L'homomorphisme (de $U$) naturel correspondant
\[
  i_B : B \longrightarrow \mathrm{Hom}_{V}(U, B)
\]
est celui qui résulte de
$B \xrightarrow{\ \sim\ } \mathrm{Hom}_{U}(U,B) \to \mathrm{Hom}_{V}(U,B)$.

Les suites $\sim$-exactes sont celles qui splittent sur $V$. Les $U$-modules
$\sim$-injectifs sont les facteurs directs des modules
$\mathrm{Hom}_{V}(U, B)$.

\page{27}
2) On a aussi
\[
  \mathrm{Hom}_{V}(A, B) \;=\; \mathrm{Hom}_{U}(U \otimes_{V} A, B) ,
\]
ce qui nous ramène à la situation \struck{duale} de celle
\uncertain{envisagée} dans les paragraphes précédents. Ainsi les objets
$\sim$-projectifs sont les facteurs directs d'objets $U \otimes_{V} A$. La
résolution canonique \ill{}

\struck{$U \otimes_{V} A \to A$}
\[
  e^{*}_{A} : A \longrightarrow U \otimes_{V} A
  \qquad \text{($V$-linéaire, donné par } a \mapsto 1 \otimes a \text{)}
\]
\[
  i^{+}_{A} : U \otimes_{V} A \longrightarrow A
  \qquad \text{($U$-linéaire surjectif, } u \otimes_{V} a \mapsto ua \text{)}
\]

3) Soit $C$ la catégorie des faisceaux de $\mathcal{O}$-modules sur $X$, et
$\widetilde{C}$ la catégorie produit $\prod_{x \in X} C^{\mathcal{O}_x}$. On a
la restriction \ill{} qui \ill{} au foncteur de Godement
$F \mapsto C^{0}(F)$.

\section*{Remarques — SGA 4 IX (pages 28 à 40)}

\page{28}
Titre porté sur la chemise : « Remarques » ; au-dessous, « SGA 4 IX ».
\note{une main d'archiviste a pointé « 11 p » en marge, ce qui est exactement
le compte des pages 30 à 40}

\page{30}
\struck{VII} IX
\begin{enumerate}
\item Sorite des faisceaux de torsion.
\item Faisceaux loc.\ constants.
\item Faisceaux constructibles.
\end{enumerate}

\uncertain{Théorie non commutative} : ce n'est \ill{} \uncertain{qui donne}
loc.\ cst.\ à fibres finies sur $X^{\mathrm{cons}}$. Le reste du n\textsuperscript{o} 2
donne les propriétés d'exactitude de la notion.

(a) \emph{Conditions équivalentes} :
\begin{enumerate}
\item[(i)] L'image inverse sur $X^{\mathrm{red}}$ est loc.\ cst.\ à fibres
  finies.
\item[(ii)] Il existe $f : X' \to X$ localement de présentation finie
  surjectif, avec $f^{*}(F)$ loc.\ cst.\ à fibres finies.
\item[(iii)] \struck{Si} $X$ quasi-compact quasi-séparé : il existe une
  partition de $X$ en $X_i$ localement fermés constructibles, en nombre fini,
  telle que $F|X_i$ soit loc.\ cst.\ à fibres finies.
\item[(iv)] \add{si $X$ quasi-compact quasi-séparé} $F$ est le noyau d'une
  double flèche $F_1 \rightrightarrows F_0$, où $F_0$ et $F_1$ sont des
  \ill{}
\item[(v)] ($X$ quasi-compact quasi-séparé) \ill{} des faisceaux représentables
  par des schémas étales de présentation finie sur $X$ ; \struck{$F$ est le
  \ill{} des constructibles de type fini}, et \ill{} d'un faisceau \ill{}
  $G$ de $F$, \ill{} $G$ \ill{} représentable, \ill{} de plus \ill{}
\item[(vi)] ($X$ quasi-compact quasi-séparé) qui admet \ill{} : $F$ admet une
  filtration \ill{} par des $i_{!}(G)$, $i : Y \to X$ immersion
  \struck{\ill{}} constructible, $G$ \ill{} sur $Y$.
\item[(vii)] ($X$ \ill{}, quasi-compact \ill{}) : $F$ admet \ill{} par
  \ill{} abélien, $A$ \ill{} ; \ill{} : $F$ \ill{} deux \ill{}
  $p_{*}\bigl( G_{X'} \bigr)$, $X'$ fini de présentation \ill{} sur $X_i$,
  \ill{} aussi.
\item[(viii)] \add{Si $A = \mathbf{Z}$, faisceaux de torsion, $X$
  quasi-compact quasi-séparé} \ill{} $F$ est réunion \ill{} de facteurs
  directs \ill{} d'un faisceau qui, \ill{}, est réunion des
  \ill{} par des faisceaux de la forme
  $p_{*}\bigl( i_{!} (\mathbf{Z}/n\mathbf{Z})_{U} \bigr)$, où
  $p : X' \to X$ est \ill{} finie et de \ill{} fini, $U$ un ouvert de $X'$ et
  $i : U \to X'$ l'immersion.
\end{enumerate}

\marginal{On peut songer aussi aux faisceaux de $A$-Modules et à la
$A$-constructibilité. Pour \ill{} de vue rédaction, il \ill{} bien de
commencer par la $\mathbf{Q}$- et la $\Delta$-constructibilité.}

(b) \emph{Déf.} Constructibles de type fini, \ill{} constructibles
\struck{et constructible \ill{}}

\emph{Cas} (c) : stabilités par $\varinjlim$ et $\varprojlim$ \ill{} (dans le
cas de la \ill{}). Il faut \struck{\ill{}} que $A$ soit noethérien.
\ill{} (cas noethérien) \ill{} par \uncertain{facteurs directs}.

\emph{Cas} (d) : constructible \ill{} (cas noethérien) \ill{} plus image
directe par \ill{}

\emph{Cas} (e) : stabilité par \struck{\ill{}} morphisme fini \ill{}
quasi-\ill{} fini \ill{}

\page{31}
\textbf{Pureté}

\textbf{Proposition.} Soit $X$ un préschéma localement noethérien
\uncertain{régulier}. \emph{Conditions équivalentes} :
\begin{enumerate}
\item[(i)] Pour \struck{tout} tout sous-préschéma régulier $Y$ de $X$, et tout
  sous-préschéma régulier fermé $Z$ de $Y$, purement de codimension $d$, on a
  \[
    H^{i}_{Z}(A_Y) = 0 \quad \text{pour } i \neq 2d
  \]
  \struck{$H^{2d}_{Z}(A_Y)$ loc.\ isomorphe à $A_Y$}
\item[(ii)] Comme \struck{dessus}, avec $d = 1$.
\item[(iii)] \struck{Comme p.\ $Y$, $Z$ comme en (i), et} Comme (i), \ill{}
  $Z$ est un diviseur de \ill{} ; on a
  \[
    H^{i}_{Z}(A_Y) = 0 \quad \text{si } i \neq 2d .
  \]
\end{enumerate}

\marginal{\ill{} de $X$}

\emph{Dém.} (i) $\Rightarrow$ (ii) trivial. (i) $\Rightarrow$ (iii).

(ii) $\Rightarrow$ (i) : par dévissage, on est ramené à \ill{}, pour $d = 1$,
$H^{2}_{Z}(A_Y)$ est loc.\ isomorphe à $A_Z$.

(iii) $\Rightarrow$ (ii) : on est ramené au \ill{} \uncertain{strict}
$X \struck{\subset} Y$ (loc.\ $2$, \ill{} $Z$ \ill{} $= 1$). À \ill{}
quelconque \ill{} (iii), que $H^{i}_{Z}(A_X) = 0$ si $i \neq 2$.
Réciproquement

\page{32}
\textbf{Prop.\ e} \quad 1\textsuperscript{o}) Faisceaux quelconques comme
$\varinjlim$ de \ill{} ; 2\textsuperscript{o}) idem d'un constructible dans
une $\varinjlim$ de constructibles ; 3\textsuperscript{o}) comparaison des
constructibles sur une $\varprojlim$ de préschémas.

\struck{Cor.\ f \quad Stabilité par image directe par un faisceau fini}

On terminera par les propriétés spéciales aux \uncertain{courbes}.

\textbf{Prop.\ f} \quad 1\textsuperscript{o}) Caractérisation de la
constructibilité par la propriété d'être loc.\ constant \ill{} ;
2\textsuperscript{o}) caractérisation loc.\ \ill{} de la clôture des
faisceaux \ill{} ; 3\textsuperscript{o}) et caractérisation des
constructibles comme objets \ill{}

\textbf{Prop.\ g} \quad Caractérisation des faisceaux loc.\ constants
\add{en $A$-Modules} \add{\ill{} constructible}

\begin{enumerate}
\item[5.] \uncertain{Variétés} \ill{} les complexes de \struck{faisceaux}
\item[4.] \struck{\ill{}} Tors et faisceaux de torsion constructibles.
\item[6.] Kummer et Artin--Schreier.
\item[7.] Dim.\ coh.\ des courbes algébriques.
\end{enumerate}
\note{une accolade portée à gauche des deux premières lignes, avec une flèche,
intervertit les n\textsuperscript{os} 5 et 4}

\emph{Leitfaden.}

\begin{tikzcd}[column sep=small, row sep=small, nodes={font=\scriptsize}]
 & 1 \arrow[d, no head] & \\
 & 2 \arrow[dl, no head] \arrow[d, no head] \arrow[dr, no head] & \\
5 & 3 \arrow[d, no head] & 6 \arrow[d, no head] \\
 & 4 & 7
\end{tikzcd}

\note{la page porte deux croquis de ce Leitfaden ; le premier, surchargé et
biffé, ne se lit pas — le diagramme ci-dessus est celui du second}

\page{33}
\struck{Théorème.} Soit $F$ un faisceau de $\ell$-torsion sur $Y$, et soit
$\xi \in H^{n}\bigl( X, f^{*}(F) \bigr)$. Pour que
$\xi \in \mathrm{Im}\, H^{n}(Y, F)$, il faut et suffit que, pour tout
$y \in Y$, la classe
$\xi_{\bar{y}} \in H^{n}(X_{\bar{y}}, F_{\bar{y}})$ induite sur la fibre
géométrique en $X_{\bar{y}}$ soit nulle.

\marginal{$Y$ loc.\ noeth.}

\textbf{Corollaire.} Supposons $f$ \uncertain{propre} \add{\ill{}} et \ill{}
acyclique par le \ill{}. Pour qu'il soit \emph{globalement} acyclique par
$\ell$, il faut et suffit que les fibres géométriques soient acycliques par
$\ell$.

\emph{Dém.\ du} \uncertain{n\textsuperscript{o} 2}. Par hypothèse
$r = 0$ :
\[
  \begin{cases}
    R^{q} f_{*}\bigl( f^{*}(F) \bigr) = 0 & \text{pour } 0 < q \leqslant n-1 \\
    R^{0} f_{*}\bigl( f^{*}(G) \bigr) \simeq G &
  \end{cases}
\]
donc la suite spectrale donne une suite exacte
\[
  \cdots \longrightarrow H^{n}(Y, F) \longrightarrow H^{n}(X, F)
  \longrightarrow H^{0}\bigl( Y, R^{n} f_{*}(f^{*}F) \bigr)
\]
\note{$\mathrm{NB.}$ on suppose $n \geqslant 0$ (ou $n \geqslant 1$ ; si
$n = 0$, cet énoncé est \ill{} un énoncé séparé avec le \ill{})}
donc on est ramené à prouver qu'une section de
$R^{n} f_{*}\bigl( f^{*}(F) \bigr)$ qui induit $0$ sur chaque fibre
géométrique est nulle. On est ramené

\page{35}
XII : 6.2.\ \ldots{} \quad
$\varphi^{i} : T^{i}(X) \to T'^{i}(X)$, \struck{\ill{}} \ill{} pour
$i \leqslant n$, \ill{} pour $i = n+1$.

$\longrightarrow A \to B \to C \to 0$

\begin{tikzcd}[column sep=small, row sep=small, nodes={font=\scriptsize}]
T^{i-1}(B) \arrow[r] \arrow[d] & T^{i-1}(C) \arrow[r] \arrow[d] & T^{i}(A) \arrow[r] \arrow[d] & T^{i}(B) \arrow[r] \arrow[d] & T^{i}(C) \arrow[d] \\
T'^{i-1}(B) \arrow[r] & T'^{i-1}(C) \arrow[r] & T'^{i}(A) \arrow[r] & T'^{i}(B) \arrow[r] & T'^{i}(C)
\end{tikzcd}
\note{le diagramme est celui que la page numérote (I) ; le quatrième terme
de la ligne du haut est d'une lecture douteuse — $T^{i}(B)$ est ce
qu'exige la suite, mais la page porte plutôt $\varphi^{i}(B)$}

a) Supposons qu'on sache que $\varphi^{i}$ est un isomorphisme pour
$i \leqslant n-1$, \ill{} pour $i = n$, et si l'on a un système de
cogénérateurs pour lesquels \ill{} $\varphi^{i}$ est un \ill{}. Alors le
diagramme (I) pour $i = n$ montre que $\varphi^{n+1}$ \ill{} un isomorphisme
pour \emph{tous} les arguments. Plus \uncertain{précisément}, si $A$ est tel
qu'il existe $A \to B$ mono et que c'est vrai pour $\varphi^{n}(B)$, c'est
vrai pour $\varphi^{n}(A)$.

b) Supposons qu'on sache que $\varphi^{i}$ est un isomorphisme pour
$i \leqslant n$, et $\varphi^{n+1}$ un mono.\ pour un système de
cogénérateurs. Alors le diagramme (I) avec $i = n+1$ montre que
$\varphi^{n+1}$ est un \ill{} pour \emph{tous} les arguments. Plus
précisément, si $\exists\, A \to B$ mono, et si c'est vrai pour $B$, c'est
vrai pour $A$.

a$'$) Même hypothèse préliminaire sur les $\varphi^{i}$ que dans a) (isom.\ si
$i \leqslant n-1$, mono si $i = n$). On sait de plus si $\varphi^{n}(B)$ est
un isomorphisme (i.e.\ une \uncertain{section}) pour $B$ donné. Si on a une
suite exacte $0 \to A \to B \to C \to 0$, il \ill{} de le vérifier pour
$\varphi^{n}(A)$ et $\varphi^{n+1}(C)$ ? Oui, \struck{parce que}
$\varphi^{n+1}(A)$ sont injectifs, \uncertain{non} en général.

b$'$) Même hypothèse préliminaire que b) ($\varphi^{i}$ un isomorphisme si
$i \leqslant n$). On veut savoir si $\varphi^{n+1}(B)$ est un mono.
Suffit-il de le vérifier pour $\varphi^{n+1}(A)$ et $\varphi^{n+1}(C)$ ?
\emph{Oui}.

\page{36}
Enfin on conclut que tout morphisme projectif \add{simple} est $\in M$, cqfd.

\textbf{Corollaire.} Supposons \struck{en \ill{}} de plus \ill{} c) la
condition suivante satisfaite :

c$'$) Soit $X' \xrightarrow{\ f\ } X \xrightarrow{\ g\ } Y$ \ill{} morphismes
\struck{\ill{}} propres, $X_0$ un sous-préschéma fermé de $X$, $U = X - X_0$,
$U' = f^{-1}(U)$. \struck{Et} Supposons que $U' \to U$ soit un isomorphisme,
\struck{et} que $X_0 \xrightarrow{\ g_0\ } Y$ \struck{soit} et
$gf : X' \to Y$ soient $\in M$. Alors $f \in M$.

Alors tout morphisme propre $f : X \to Y$ \struck{est} $\in M$.

En effet, il est trivial que c$'$) implique c), donc tout morphisme projectif
est $\in M$. On obtient le cas d'un morphisme \struck{quelconque} propre
quelconque en utilisant c), \ill{} qui provient du \ill{} le
\uncertain{dévissage} \ill{} \struck{\ill{} c$'$} et d'utiliser \ill{}
\uncertain{noethérienne}. Le lemme \struck{\ill{}} de Chow \ill{} donne alors
le résultat voulu.

\page{37}
c) Supposons les $\varphi^{i}$ \struck{des isomorphismes} \struck{fonctoriels}
\add{un épimorphisme surjectif}. Quand peut-on conclure que $\varphi^{n+1}$
est un mono ? Réponse : il \uncertain{suffit} que $T'^{n+1}$ soit effaçable.
Donc si les $T'^{i}$ sont effaçables pour $i > i_0$, et bien qu'il en donne
\ill{} $\leqslant i_0$, on \ill{}

c$'$) \struck{Supposons les $\varphi^{i}$ des isomorphismes \ill{} ; on n'a
pas \ill{} $T'^{i-1}$ \ill{} sont effaçables} \ill{} pour $i > i_0$, et
$\varphi^{i}$ un épi.\ \ill{} $i \leqslant n$ (ou $n \geqslant i_0$). Alors
$\varphi^{i}$ est un \struck{épi} \emph{isom} pour $i \leqslant n$, un épi en
degré $n+1$.

Revenons aux a), b). On a des propriétés \struck{\ill{}} $I_n$, $J_n$ d'un
argument $A \in \uncertain{\mathrm{Ob}\, C}$ :
\begin{quote}
$I_n(A)$ : $\varphi^{i}(A)$ isom.\ si $i \leqslant n$ ;\\
$J_n(A)$ : $\varphi^{i}(A)$ isom.\ si $i < n$, mono si
\uncertain{$i = n+1$}.
\end{quote}
\begin{quote}
$I_n \Rightarrow J_n \Rightarrow J_m \cdots$\\
$I_n$ pour tout $A$, $+$ $J_n$ pour \ill{} $\Rightarrow$ $J_n$ pour tout $A$\\
$J_n$ pour tout $A$, $+$ $I_{n+1}$ pour tout \ill{} $\Rightarrow$ $I_{n+1}$
pour tout $A$
\end{quote}

\emph{Concl.} Pour vérifier $I_n$ (ou $J_n$) pour tout $A$, sachant que c'est
vrai pour \ill{} \uncertain{fixés}, il suffit de vérifier $I_m(B)$ (ou
$J_m(B)$) pour $m \leqslant n$ et tout \ill{} $B$.

On avait ici des arguments \ill{} \ill{}. Dans les questions de dim.\ coh.,
\ill{} il faut « démouler » la question et des \ill{} des conditions pour que
l'un $T^{i}$ soit nul (pour tout argument). Ceci en \ill{} utiliser les faits
suivants (savoir aussi dans \ill{} que $T^{i}$ est
\note{la suite de cette page se lit page 39}

\page{38}
\struck{Lemme 1} \textbf{Proposition.} Soit $M$ un ensemble de morphismes de
préschémas loc.\ noethériens. On suppose que
\begin{enumerate}
\item[a)] pour un morphisme $f : X \to Y$, l'\ill{}-$M$ \ill{} : $M$ est local
  sur $Y$ ;
\item[d)] si $X \xrightarrow{\ f\ } Y \xrightarrow{\ g\ } Z$ sont tels que $f$
  soit une immersion fermée et $gf \in M$ \add{et $gf$ simple}, alors
  $gf \in M$ ; \add{$f : \mathbf{P}^{i}_{Y} \to Y$ est \ill{} ; et
  $\mathrm{id}_Y : Y \to Y$ est $\in M$} ;
\item[b)] si $X \xrightarrow{\ f\ } Y \xrightarrow{\ g\ } Z$ avec
  $f, g \in M$, alors $gf \in M$ ;
\item[c)] si $X = \mathbf{P}^{r}_{Y} \xrightarrow{\ g\ } Y$ est le morphisme
  structural canonique, \struck{et si}
  $X' = \mathbf{P}^{1}_{Y} \times \mathbf{P}^{r-1}_{Y}$ est muni du morphisme
  birationnel bien connu $X' \xrightarrow{\ f\ } X$, et si $gf \in M$, alors
  $g \in M$.
\end{enumerate}

Sous ces conditions, tout morphisme \add{simple} projectif $f : X \to Y$ est
$\in M$.

En effet, on voit par récurrence sur $r$, utilisant a), b), c), que les
$\mathbf{P}^{r}_{Y} \xrightarrow{\ f\ } Y$ sont $\in M$. Utilisant d), il
s'ensuit que pour tout sous-préschéma fermé $X$ de $\mathbf{P}^{r}_{Y}$
\add{simple$/Y$} la projection sur $Y$ est $\in M$.

\page{39}
Deux énoncés :
\begin{quote}
si $T^{i}(B) = 0$, alors $T^{i}(A) = 0$ pour tout facteur direct $A$ de $B$ ;\\
si $A$ a une filtration \uncertain{par} \ill{} facteurs $A_{\alpha}$ telle que
$T^{i}(A_{\alpha}) = 0$ pour tout $\alpha$, alors $T^{i}(A) = 0$.
\end{quote}

Pour les foncteurs cohomologiques de faisceaux de $\ell$-torsion
\emph{constructibles} sur $X$ (\ill{}), \ill{} remarquer \struck{déjà}
\ill{} que, pour avoir $T^{i}(F) = 0$ pour tout tel $F$, il suffit de le
prouver pour un $F$ de la forme
\[
  i_{!}\Bigl( p_{*}\bigl( (\mathbf{Z}/n\mathbf{Z})_{Y'} \bigr) \Bigr)
  \;=\; q_{*}\, j_{!}\bigl( (\mathbf{Z}/n\mathbf{Z})_{Y'} \bigr) ,
\]
$Y$ étant un sous-préschéma constructible de $X$, $Y'$ un revêtement étale
fini de $Y$ de degré constant, et $p : Y' \to Y$, $i : Y \to X$ les morphismes
canoniques. On les prend d'un \ill{} immédiatement des $X'$ fini sur $X$.

\begin{tikzcd}
Y' \arrow[r, hook] \arrow[d, "p"'] & X' \arrow[d, "q"] \\
Y \arrow[r, hook, "i"'] & X
\end{tikzcd}

Toute sous-catégorie pleine de la catégorie des faisceaux de $\ell$-torsion
constructibles sur $X$, stable par facteurs directs, par extension, et
contenant les faisceaux
$q_{*}\, j_{!}\bigl( (\mathbf{Z}/n\mathbf{Z})_{Y'} \bigr)$, \ill{}, est
\ill{}. Si l'on veut une sous-catégorie \struck{\ill{}} telle que, pour toute
suite exacte $0 \to A \to B \to C \to 0$, si $B$ et $C$ y sont, alors $A$ y
est — je crois si c'est une sous-catégorie \emph{épaisse} — alors il suffit
de vérifier que les $q_{*}\bigl( (\mathbf{Z}/n\mathbf{Z})_{X'} \bigr)$ y
sont, $X'$ fini sur $X$.

Car $X$ \ill{}, on peut \add{de plus} prendre $X'$ intègre. Si $X$ est
\struck{\ill{}} \add{de plus} \ill{} noeth.\ séparé, on peut prendre $X'$
intègre normalisé.

\page{40}
\note{la page entière est biffée en croix}

\struck{sont} les \ill{} de $S$, i.e.\ dans $X_0$. Donc, pour voir que c'est
un isomorphisme \ill{} de $X_0$, il suffit de voir que l'\ill{} \ill{}
$R\, j^{!}$, où $g : X_0 \to Y$, \ill{} \ill{} un isomorphisme en $2$. On
\ill{}

\[
  R\, g^{!}\bigl( R\, i^{!}(\mu_{\alpha}) \bigr) \longrightarrow
  R\, j^{!}\bigl( A_Y[-2] \bigr)
\]
\note{la page suspend chacun des deux termes, par un signe $\wr$, à une
seconde ligne qui ne se lit pas}

\marginal{$X_0 \xrightarrow{\ g\ } Y \xrightarrow{\ j\ } X$, la composée
notée $h$}

\struck{$R\, h^{!}(\mu_{\alpha})$}

\end{document}
